\chapter{缓存、TLB、预取与数据布局}

CPU 执行单元的速度很快，但它只能处理已经到达寄存器的数据。数据从内存到寄存器之间要经过多级缓存、TLB、加载队列、存储队列和预取器。AI 算子的大部分优化都在处理同一个矛盾：数学上相邻的元素，不一定在内存里相邻；代码里重复使用的值，不一定能在缓存里留到下一次使用。

\section{缓存层级不是抽象背景}

缓存以 cache line 为基本单位搬运数据，常见大小是 64 字节。读取一个 float 时，硬件通常会把同一条 cache line 里的相邻元素一起带进来。如果后续代码顺序访问这些元素，带宽利用率高；如果每次跨很大步长访问，只用到每条 cache line 里的一个元素，实际带宽会被浪费。

矩阵按行主序存储时，连续访问一行是顺序访问，连续访问一列是跨步访问。第一册已经讲过地址计算，第二册要进一步问：这个地址序列对 cache line 是否友好，对硬件预取器是否友好，对 TLB 是否友好。一个三重循环的数学结果不变，但循环顺序不同，内存访问轨迹完全不同。

\begin{lstlisting}[language=C++,caption={同一个矩阵，行访问和列访问的内存轨迹不同}]
#include <cstddef>
#include <cstdint>
#include <span>

float sum_by_rows(std::span<const float> a,
                  std::int32_t rows,
                  std::int32_t cols) {
    float acc = 0.0F;
    for (std::int32_t r = 0; r < rows; ++r) {
        for (std::int32_t c = 0; c < cols; ++c) {
            acc += a[static_cast<std::size_t>(r) * cols + c];
        }
    }
    return acc;
}

float sum_by_cols(std::span<const float> a,
                  std::int32_t rows,
                  std::int32_t cols) {
    float acc = 0.0F;
    for (std::int32_t c = 0; c < cols; ++c) {
        for (std::int32_t r = 0; r < rows; ++r) {
            acc += a[static_cast<std::size_t>(r) * cols + c];
        }
    }
    return acc;
}
\end{lstlisting}

当矩阵很小，两个函数都在缓存里跑，差距可能不明显；当矩阵超过缓存容量，列访问就会暴露跨步访问的问题。AI 计算中的很多布局变换，例如 NHWC 和 NCHW 的选择、权重预打包、KV Cache 的排布，最终都要回到这个问题：下一次访问的数据，是否已经在离核心足够近的层级里。

\section{TLB 和页粒度}

缓存按 cache line 管数据，TLB 按页缓存虚拟地址到物理地址的翻译。访问大量离散页时，即使数据总量不算夸张，也可能因为 TLB miss 变慢。大模型推理中，权重很大，KV Cache 随序列增长，内存访问不仅要考虑缓存容量，还要考虑页表翻译和 NUMA 放置。

TLB 问题常被初学者忽略，因为代码里看不到它。一个指针加下标的表达式背后，硬件要完成地址生成、权限检查和地址翻译。若访问模式顺序，TLB 和预取器都容易工作；若访问模式随机，硬件很难提前准备。Embedding lookup、稀疏访问、哈希结构和不规则图计算通常比规则矩阵乘更容易受 TLB 和缓存 miss 影响。

\section{AoS、SoA 与结构化数据}

\term{AoS}{Array of Structures} 把一个对象的多个字段放在一起，适合按对象访问；\term{SoA}{Structure of Arrays} 把每个字段分别放成数组，适合 SIMD 和批量处理。AI 算子通常处理大批同类型数值，因此 SoA 或接近 SoA 的布局更常见。系统代码中也会遇到两者取舍：一个推理请求对象可能适合 AoS，但一批 token 的 logits 或概率更适合连续数组。

\begin{lstlisting}[language=C++,caption={AoS 和 SoA 的访问形状}]
#include <cstdint>
#include <vector>

struct TokenStateAos {
    float logit;
    float probability;
    std::int32_t token_id;
};

struct TokenStateSoa {
    std::vector<float> logits;
    std::vector<float> probabilities;
    std::vector<std::int32_t> token_ids;
};
\end{lstlisting}

如果代码要对所有 logit 做 Softmax，SoA 更容易顺序扫描，也更容易向量化。如果代码经常拿某一个 token 的全部字段做复杂逻辑，AoS 可能更自然。布局没有绝对好坏，只有访问模式是否匹配。

\section{预取与局部性边界}

硬件预取器擅长发现顺序和固定步长访问，不擅长预测复杂间接访问。手动预取有时有用，但必须谨慎。预取太晚，数据仍然来不及；预取太早，数据可能在使用前被挤出缓存；预取太多，会污染缓存并占用带宽。本册会优先强调通过布局和循环顺序创造自然局部性，而不是一开始就依赖手动预取。

\begin{keyidea}
缓存优化不是把所有数据都塞进缓存，而是让即将重复使用的数据在重复使用前不要离开合适的缓存层级。分块、packing、循环交换和算子融合都是围绕这个目标展开。
\end{keyidea}

\begin{exercisebox}
把一个大矩阵分别按行和按列求和，记录不同矩阵规模下的运行时间。矩阵从 L1 cache 能容纳的规模开始，逐步扩大到 L2、L3 之外。观察差距从哪里开始变大，并结合 cache line 解释原因。
\end{exercisebox}
